\documentclass[7pt]{article}
%\usetheme{Warsaw}

%\usepackage{beamerinnerthemecircles, beamerouterthemeshadow}
%\usepackage{scrextend}

%\changefontsizes{7pt}


% Standard packages
\usepackage{graphicx}
\usepackage{graphics}

\usepackage{color}
\usepackage[OT4]{fontenc}
\usepackage[utf8]{inputenc}
%\usepackage[latin1]{inputenc}
\usepackage{times}\usepackage{amsmath}
\usepackage[T1]{fontenc}

%\mode<presentation> {
%    \usetheme{Singapore} %Bergen, Berkely, Berlin, Boadilla, CambridgeUS, Darmstadt,
%                          %Frankfurt, Goettingen, Singapore, Warsaw
%    \usecolortheme{default} %beetle, seahorse, wolverine, dolphin, beaver
%}


\title{2D membrane vibrations (FD version)}

\definecolor{whi}{rgb}{0.95,0.95,0.95}
% Setup TikZ

%\usepackage{tikz}
%\usetikzlibrary{arrows}
%\tikzstyle{block}=[draw opacity=0.7,line width=1.4cm]


% Author, Title, etc.




% The main document

\begin{document}
\maketitle




\section{the problem}


We solve the wave equation for $\Psi$ its deviation from the equlibrium position at ($x,y$) point
with damping 
\begin{equation} \frac{\partial^2 \Psi}{\partial t^2}=v^2\nabla^2\Psi -2d\frac{\partial u}{\partial t}\end{equation}
for a square membrane $(x,y)\in(0,L)\times(0,L)$ using the finite element method. We set the speed of sound equal to $v=1$, the size $L=1$,
and the damping constant $2d=1$. In the central node of the mesh $(c,c)$ we apply a periodic force $\sin\omega t$

We use a finite difference approach on a 2D mesh. 
Upon discretization of the derivatives with the usual approach for the second one and $\frac{\partial u}{\partial t}\simeq \frac{u(t)-u(t-\Delta t)}{\Delta t}$
we get 
\begin{eqnarray}
u_{ij}(t+\Delta t)=u_{ij}(t)[2-2d \Delta t]+u_{ij}(t-\Delta t)(2d \Delta t -1)+\\\frac{\Delta t^2 v^2}{\Delta x^2}(u_{i+1,j}(t)+u_{i-1,j}(t)+u_{i,j+1}(t)+u_{i,j-1}(t)-4u_{i,j}(t)) \\
+\delta_{i,c}\delta_{j,c}\Delta t^2\sin(\omega t).
\end{eqnarray}

We start from a standing membrane $u(0)=u(\Delta t)=0$. 
$\Delta x=L/N$, for $N\times N$ mesh. Lets set $v\Delta t=\Delta x/2$.
We monitor the energy $E(t)=0.5\int_0^Ldx\int_0^Ldy\left[(\frac{\partial u}{\partial t})^2+(\frac{\partial u}{\partial x})^2+(\frac{\partial u}{\partial y})^2\right]$.
We integrate the equation till $t=80$. 

\section{Deliverables}
\begin{itemize} \item
Lets set $\omega=\pi$. Does the motion obtained a periodic form?
\item 
Lets us assume that the first 10 units of time is necessary for the membrane vibrations to reach their stationary form.
We evaluate the average energy over the time interval $(10,80)$.
The average energy is $\langle E \rangle(\omega)=\frac{1}{70} \int_{10}^{80} dt E(t)$.

Plot $\langle E \rangle(\omega)$ for $\omega=(1.5,15)$.
\item Determine the resonant frequencies.
\item How the resonant peaks react to the change of the damping constant?
\end{itemize}
\end{document}




